\scanpage{236}
\mspage{17}
\margin{n\textsuperscript{o} 227}

des points fixes.

PROPOSITION 1. - \uline{Soit $\mathcal{L}$ un système linéaire sans composante fixe, et $T$ l'application rationnelle associée. Soit $P'$ un point normal de $V$. Alors $T$ est définie en $P'$ si et seulement si $P'$ n'est pas un point de base.}

\uline{Dém.} Supposons que $P'$ ne soit pas un point de base. Soit $Y=\sup$ des pôles des $f_i$, et écrivons $(f_i)=X_i-Y$. Alors par hypothèse, il existe un index $i$, mettons $i=0$, tel que $P'\notin|X_0|$. On a $(f_i/f_0)=X_i-X_0$. Par la prop. 5 du \S 1, chaque fonction $f_i/f_0$ n'a pas de pôle passant par $P'$, et est définie en $P'$. La réciproque est également claire.

Soit $\mathcal{L}$ un système linéaire sans composante fixe, et $T:V\longrightarrow V'$ l'application rationnelle induite par $\mathcal{L}$. On va voir que $T$ applique les diviseurs de $\mathcal{L}$ sur les \inserted{\unsure{sections}} \unsure{hyperplans|hyperplanes} de $V'$ de la manière suivante. Prenons $X\in\mathcal{L}$, $P'\in|X|$ et supposons que $P'$ ne soit pas un point de base pour $\mathcal{L}$. Prenons un diviseur $X_0\in\mathcal{L}$ tel que $P'\notin|X_0|$. Soit $L_0$ l'espace dérivé de $\mathcal{L}$ par $X_0$ et $f_0,\ldots,f_n$ une base de $L_0$. Alors chaque fonction $f\in L_0$ peut être écrite
\[
f=f_c=\sum c_i f_i\qquad c_i\in k
\]
et $(f_c)=X_c-X_0$. Supposons $X=X_c$. Puisque $P'\notin|X_0|$, toutes les fonctions $f_i$ sont définies en $P'$. Mais $P'\in|X_c|$ implique que $\sum c_i f_i(P')=0$, i.e. $T(P')$ est dans l'hyperplan définit par l'équation $\sum c_i X_i=0$.

Un système linéaire sera dit \uline{ample} \inserted{s'il est sans \unsure{com.|comp.} fixe et} si l'application rationnelle $T$ induite par $\mathcal{L}$ est un isomorphisme de $V$ sur $V'$. Si $D$ est un diviseur tel que $\mathcal{L}(D)$ soit ample, alors $D$ est dit \uline{ample}. Le système des

\mspage{235}
